\bookchapter{Capítulo 32: Manual de emergência de segurança do celular}

\chapterintro{Este capítulo é um “plano de emergência”. Espero que você nunca precise, mas se o celular for perdido, roubado ou se você cair em golpe, este manual ajuda a minimizar os prejuízos. Em desastres como terremotos e enchentes, o celular também é ferramenta para receber alertas, pedir socorro com localização e manter contato. Leia agora, salve capturas dos passos-chave ou imprima e deixe em casa — na hora do ocorrido o aparelho pode não estar com você, e em desastres pode faltar bateria ou sinal.}

\booksection{Cenário zero: terremoto / clima extremo — como o celular ajuda a salvar sua vida}

Aqui o foco é a segurança física, não contas e dinheiro. Configurações de alerta no Capítulo 13 «Alertas climáticos, de desastres naturais e de terremotos»; aqui falamos do que fazer depois que o alarme toca.

\textbf{Alarme de terremoto: abrigue-se primeiro, saia depois — não fique só gravando vídeo}\par

\begin{enumerate}
  \item Abrigue-se perto: abaixe, cubra e segure-se (sob a mesa, canto de parede estrutural, móvel firme), proteja cabeça e pescoço; afaste-se de vidros, paredes externas e lustres.
  \item Não corra para as escadas: o tremor principal costuma durar dezenas de segundos; as escadas podem ser mais perigosas. Nunca use o elevador.
  \item Depois que o tremor passar, evacue com calma: use as escadas, atenção a réplicas e objetos que caem; vá para área aberta, longe de postes, placas e construções precárias.
  \item Use o celular com sabedoria: mensagem curta para contatos de emergência ("Estou em XX, seguro/ferido, preciso de ajuda") consome menos bateria que uma ligação longa; se der, prefira texto (SMS/WeChat).
  \item Sem sinal: veja o Capítulo 7 (SOS por satélite / mensagens via satélite, se o aparelho suportar); ou tente em local alto e aberto; em alguns casos, o rádio walkie-talkie por Bluetooth sem internet também ajuda.
  \item Desligue o desnecessário: em emergência, ative o modo avião e abra só no horário preciso; desligue varredura de Bluetooth e reduza o brilho — ganha algumas horas. Guarde a última barra de bateria para avisar que está bem e pedir socorro.
\end{enumerate}

\begin{tcbwarn}{Chuva forte / tufão / enxurrada: aja já na fase de alerta}
\begin{itemize}
  \item Ao ver alerta vermelho: reduza a saída, evite espaços subterrâneos, viadutos, margens de rio e encostas.
  \item Se ficar preso na rua: suba para andares altos de construções firmes; não atravesse alagamentos (podem faltar tampas de bueiro e haver fuga de energia); ao ligar 119/110, informe o local e o nível da água.
  \item Antes de sair: confira o risco no app de clima e nas contas oficiais de alerta; em viagens longas de carro, evite trechos de montanha junto a rios e barrancos.
  \item Combinado em família: "se perdermos contato, vale a última localização no WeChat/SMS", para não gastar bateria ligando loucamente.
\end{itemize}

\end{tcbwarn}

\begin{tcbpractice}{4 coisas para fazer desde já}
\begin{enumerate}
  \item Ative o alerta de terremoto do sistema ou um app confiável e permita que ele ignore o silencioso.
  \item Fixe 2 a 3 "contatos de emergência" no topo da agenda e exiba-os na tela de bloqueio (cartão médico iOS/Android ou informações de bloqueio).
  \item Tenha power bank em casa e na bolsa; antes de viagens longas, carregue tudo e leve cabo.
  \item Memorize ao menos um número que dê para ligar sem a agenda (o celular de um familiar) e anote no papel, na carteira.
\end{enumerate}

\end{tcbpractice}

\booksection{Cenário um: celular perdido/roubado — os 30 minutos de ouro}

Os primeiros 30 minutos são “tempo de ouro”: com ação certa, você protege dinheiro e privacidade. Siga a ordem abaixo, sem pular etapas:

\textbf{Procedimento padrão nos 30 minutos seguintes}\par

\begin{enumerate}
  \item Ligue para o seu número com o celular de outra pessoa (minuto 1): talvez tenha caído perto e alguém de bem atenda.
  \item Bloqueie o aparelho remotamente (minutos 2–5): use outro celular ou o computador para entrar no serviço de localização: | | iPhone: iCloud.com > Buscar > escolha o aparelho > Marcar como perdido
  \item Android por marca: Huawei cloud.huawei.com / Xiaomi i.mi.com / OPPO cloud.oppo.com etc.
  \item Bloqueie o SIM (minutos 5–10): ligue para a operadora (China Mobile 10086 / Unicom 10010 / Telecom 10000) e peça a suspensão do SIM. É o passo mais crítico — quem tiver o seu chip recebe os códigos e redefine todas as senhas.
  \item Congele pagamentos (minutos 10–20): | | Alipay: 95188 ou, no celular de outra pessoa, Alipay > Eu > Configurações > Central de Segurança > Congelamento de emergência
  \item WeChat: registre o boletim e fale com o suporte do WeChat, ou use Central de Segurança do WeChat > Congelar conta
  \item App do banco: ligue para o número no verso do cartão e peça o congelamento do internet banking
  \item Avise familiares e amigos (minutos 20–25): publique nos Moments ou peça a alguém da família para avisar que o celular sumiu, evitando golpes com a sua identidade.
  \item Registre boletim (minutos 25–30): ligue 110 ou vá à delegacia mais próxima e peça o comprovante. Ele será necessário para seguros e disputas bancárias.
\end{enumerate}

\booksection{Cenário dois: você caiu em golpe — janela de 72 horas}

Descobriu que foi golpe? O mais importante é manter a calma. Decisões no pânico causam segunda perda. Siga estes passos:

\begin{enumerate}
  \item Pare imediatamente qualquer transferência: não importa o que digam, não envie mais um centavo. "Mande mais uma taxa de garantia que devolvemos o resto" é o roteiro clássico do golpe do golpe.
  \item Guarde todas as provas: não apague conversas nem histórico de chamadas! Tire capturas de chats, transferências e dados da conta do golpista. São as provas-chave para o boletim e a recuperação do dinheiro.
  \item Ligue 96110 agora: é a linha nacional antifraude. O atendente orienta o bloqueio emergencial — se a transferência foi recente (em geral até 24h), o banco pode congelar a conta do golpista.
  \item Ligue 110: leve celular, documento e comprovantes à delegacia para prestar depoimento.
  \item Fale com o banco/plataforma de pagamento: relate o golpe e peça o congelamento das contas envolvidas. Se foi fraude no cartão de crédito, solicite o chargeback.
  \item Troque todas as senhas: em outro aparelho seguro, redefina Apple ID/conta da marca, WeChat, Alipay e apps bancários, e ative a autenticação em duas etapas.
\end{enumerate}

\booksection{Cenário três: vírus/cavalo de Troia no celular — o que fazer?}

Se aparecerem anúncios misteriosos, apps estranhos instalados sozinhos, conta de telefone nas alturas ou consumo anormal de dados, provavelmente há malware:

\begin{itemize}
  \item Passo 1: modo de segurança — no Android, desligue segurando o botão de energia, ligue novamente e, ao vibrar, segure o volume para baixo até entrar no modo de segurança (aparece "Modo de segurança" no canto inferior esquerdo). Assim, apps de terceiros ficam desativados e fica fácil investigar.
  \item Passo 2: desinstale apps suspeitos — em Configurações > Apps, ordene por data e remova os apps desconhecidos instalados recentemente.
  \item Passo 3: escaneie com um app de segurança — instale o "Tencent Mobile Manager" ou o "360 Mobile Guard" e faça uma varredura completa.
  \item Passo 4: restaurar as configurações de fábrica (última instância) — se nada funcionar, faça backup dos dados importantes e restaure. Atenção: tudo será apagado; salve fotos, contatos e arquivos antes.
  \item O iPhone praticamente não pega vírus — o sandbox do iOS dificulta a sobrevivência de malware. Se o iPhone agir estranho, o mais provável é bug de algum app ou configuração, não vírus. A restauração resolve 99\% dos casos.
\end{itemize}

\booksection{Cenário quatro: conta invadida — recupere o controle rapidamente}

\begin{enumerate}
  \item Troque a senha na hora: se ainda conseguir entrar, redefina imediatamente. A nova senha deve ser completamente diferente da antiga.
  \item Desconecte outros aparelhos: em Segurança da conta > Gerenciar dispositivos, derrube tudo o que não reconhecer.
  \item Confira o histórico de acessos: veja os últimos 7 dias; login de outra cidade confirma a invasão.
  \item Ative a autenticação em duas etapas: se ainda não tem, ative agora. É o método mais eficaz para não ser invadido de novo.
  \item Verifique os dados vinculados: telefone e e-mail da conta. Se foram trocados, use a função de recuperação/apelação.
  \item Confira os pagamentos: em WeChat/Alipay, veja se pagamentos com um toque ou débitos automáticos foram alterados e se adicionaram cartões a contas estranhas.
\end{enumerate}

\booksection{Cenário cinco: mensagem de chantagem — não tenha medo, não pague}

Dois golpes comuns: (1)“tenho um vídeo/foto sua, pague ou envio para todos os seus contatos”; (2)“seu celular foi hackeado, pague ou travo tudo”. Nesse caso:

\begin{itemize}
  \item 99,9\% é golpe: eles não têm nada seu. Enviaram a mesma chantagem para dezenas de milhares; sempre alguém se assusta e paga.
  \item Não responda e não pague: responder mostra que "alguém usa este número", e você vira alvo fácil de mais assédio.
  \item Tire captura e guarde as provas; depois, bloqueie o contato.
  \item Troque a senha: por precaução, redefina as senhas das contas principais.
  \item Se houver risco real de vazamento: registre boletim. A polícia tem experiência nesses casos.
\end{itemize}

\booksection{Cenário seis: golpe com deepfake de voz/rosto — ver não é mais crer}

Desde 2024, a IA consegue “transformar” o golpista em um familiar, chefe ou amigo no vídeo. “Mãe, tive um problema, transfere logo” — a pessoa parece e soa igualzinha ao seu filho, mas é falsa. É um dos golpes mais difíceis e perigosos de hoje.

\textbf{Passo 1: identifique falhas — quatro brechas comuns do deepfake}\par

\begin{itemize}
  \item Piscadas fora do normal: pessoas piscam 15–20 vezes por minuto; a IA costuma piscar pouco ou com ritmo rígido.
  \item Lábios fora de sincronia: há frações de segundo entre a boca e o som.
  \item Contorno borrado: orelhas, cabelo e bordas do rosto podem tremer ou ficar irregulares.
  \item Luz inconsistente: os realces e sombras do rosto não batem com o fundo; de perfil, os traços "giram" de forma estranha.
\end{itemize}

\begin{tcbinfo}{Passo 2: três testes para saber se é real}
\begin{itemize}
  \item Peça ações específicas: vire a cabeça, cubra o rosto, tape a boca e abra/feche a mão. O deepfake costuma "entregar" em oclusões e giros amplos — o rosto deforma ou não acompanha.
  \item Faça uma pergunta privada: invente algo como "em que hospital levamos o cachorro da última vez?" — o golpista não saberá ou gaguejará.
  \item Ligue de volta para o número que você já conhece: desligue o vídeo e disc o contato salvo. O golpista não controla o número real; na hora da ligação, a farsa cai.
\end{itemize}

\end{tcbinfo}

\begin{tcbwarn}{Passo 3: o que fazer se confirmar o golpe}
\begin{enumerate}
  \item Pare transferências: por mais "urgente" que pareça, pare. Quem insiste em "transfere já, não desliga" é quase sempre golpista.
  \item Guarde todas as provas: grave a tela, capture o vídeo, conversas, conta e comprovantes — são materiais cruciais para a perícia.
  \item Ligue 96110: informe que se trata de golpe com IA (deepfake) e o atendente iniciará o bloqueio emergencial.
  \item Ligue 110: leve as provas à delegacia e destaque o uso de IA falsificada, facilitando a classificação do caso.
  \item Avise família e amigos: o golpista pode usar o seu vídeo falso para enganar outras pessoas. Publique um alerta no Moments ou no grupo da família.
\end{enumerate}

\end{tcbwarn}

\booksection{Cenário sete: ataque de SIM swap — seu número foi “roubado”}

Pouca gente ouviu falar, mas o estrago é enorme. No SIM swap, o atacante falsifica seus dados e pede um chip novo na loja da operadora em seu nome. Assim que o chip novo ativa, o seu sinal some e o atacante passa a receber todos os códigos por SMS — Alipay, banco, WeChat, tudo aberto.

\begin{tcbinfo}{Como perceber? Três sinais de alerta}
\begin{itemize}
  \item O celular de repente mostra "Sem serviço"/"Sem SIM", mesmo sem você mexer no chip e sem resolver ao reiniciar — costuma ser o momento do sequestro.
  \item Chega SMS da operadora confirmando "reposição de SIM" ou "serviço ativado" — mas você não pediu nada.
  \item Deixa de receber códigos, e também outros SMS — as mensagens foram redirecionadas para o chip do atacante.
\end{itemize}

\end{tcbinfo}

\begin{tcbwarn}{Proteção — coloque senha no SIM}
\begin{itemize}
  \item Defina o PIN do SIM: Configurações > Telefone > PIN do cartão SIM. Ativado, cada reinício ou troca de chip pede o PIN. Mesmo que peguem o cartão, não usam.
  \item Ative a verificação da operadora: nos apps China Mobile/Unicom/Telecom, procure "Central de Segurança" ou "Verificação de serviços" e ligue o alerta de reposição de SIM e a senha extra. Assim, quem for atrás do seu chip passará por checagem adicional.
  \item Contas importantes não só por SMS: em banco, Apple ID etc., prefira autenticador (Google Authenticator, Microsoft Authenticator) ou chave física ao código por SMS.
  \item Cópias de documento com marca d'água: escreva "uso exclusivo para XX" — atacantes costumam se passar por você com dados de documento vazados.
\end{itemize}

\end{tcbwarn}

\textbf{Se o SIM for sequestrado — velocidade é tudo}\par

\begin{enumerate}
  \item Ligue imediatamente para a operadora e peça o bloqueio: China Mobile 10086 / Unicom 10010 / Telecom 10000. Solicite suspensão e congelamento da reposição para cortar o uso do seu número.
  \item Entre nas contas por outro aparelho: use o PC ou o celular de um familiar para trocar senhas de banco, Alipay, WeChat e Apple ID e desconectar dispositivos estranhos.
  \item Congele pagamentos: Alipay 95188, WeChat Central de Segurança, e o banco para o internet banking.
  \item Vá à loja da operadora com documento: peça um chip novo presencialmente e recupere o número.
  \item Registre boletim: o SIM swap configura crime contra dados pessoais; leve provas (horário da suspensão, prejuízo financeiro, registros de reposição).
  \item Auditoria completa: depois de recuperar o número, revise acessos recentes e extratos para ver se houve uso indevido.
\end{enumerate}

\begin{tcbwarn}{Prevenir é melhor — 5 hábitos diários}
\begin{enumerate}
  \item Senha de tela de bloqueio + autenticação em duas etapas, os dois ativos
  \item Confirme que "Localizar meu celular" está ativado
  \item Defina o PIN do SIM (Configurações > Telefone > PIN do SIM); se perder o celular, não usam o seu chip
  \item Não instale apps de origem duvidosa nem clique em links estranhos
  \item Faça backup periódico de fotos, contatos e conversas
\end{enumerate}

\end{tcbwarn}

\begin{tcbpractice}{Prática: crie o seu "cartão de emergência"}
\begin{enumerate}
  \item Em um papel, anote: operadora (China Mobile 10086 / Unicom 10010 / Telecom 10000), banco, 96110 antifraude, 110 e telefones da família.
  \item Anote o e-mail de cadastro do Apple ID/conta da marca e um contato alternativo.
  \item Guarde o papel na carteira ou cole na geladeira — não deixe só no celular, porque na emergência ele pode não estar com você.
  \item Confira agora: "Localizar meu celular" ativado e PIN do SIM definido?
\end{enumerate}

\end{tcbpractice}

\begin{tcbsummary}{Resumo do capítulo}
\begin{itemize}
  \item Celular perdido — 30 minutos de ouro: ligar do celular de alguém bloquear remotamente suspender SIM congelar pagamentos avisar contatos boletim
  \item Após golpe — 72 horas: parar transferências guardar provas 96110 110 banco trocar senhas
  \item Malware: modo de segurança remover apps suspeitos varredura restaurar fábrica (última instância)
  \item Conta invadida: trocar senha derrubar dispositivos ver acessos ativar 2FA checar vínculos checar pagamentos
  \item Chantagem é 99,9\% golpe: não responda, não pague, capture e bloqueie
  \item Deepfake: olhe piscadas/lábios/bordas/luz; peça para virar e cobrir o rosto, faça pergunta privada e ligue de volta
  \item SIM swap: "Sem serviço" do nada desconfie; PIN do SIM, verificação da operadora e autenticador nas contas importantes
  \item 5 hábitos: bloqueio + 2FA, localizar celular, PIN do SIM, sem apps duvidosos, backup regular
\end{itemize}

\end{tcbsummary}
